% chapters/12_conclusions.tex -- writer R3 (W6).
\chapter{Conclusions}\label{ch:conclusions}

This chapter states what the study has established about the thickness dependence of the three single IWO bottom-gate TFTs (2, 6.3 and 13.2~nm) and about the ATLAS model built to describe them. The results are ordered by strength of evidence, exactly as in \file{analysis_2026-09-25/SYNTHESIS.md}. Demonstrated conclusions rest on measured curves and converged simulations. Strongly supported interpretations rest on several independent lines of evidence but are not uniquely proven. Plausible hypotheses have not been verified, and rejected explanations are listed together with the quantitative inconsistency that rejects them. Every statement about data concerns one device per thickness ($N = 1$). The chapter ends with the strongest claims that the study can support in a publication, the corrections that must be applied to earlier documents of the package, and an outlook.

\begin{statusbox}{Overall status}
The model is calibrated, not validated. The out-of-sample test of the shared parameters within the data (6.3~nm threshold, \run{0014}) failed by $-0.294$~V, and the registered hypothesis tests B0 and C1 for the 6.3~nm film failed or partly failed. The ID-VD and 338/358~K predictions are registered with hashes in two exact mobility variants and await independent data (\cref{ch:predictions}).

Model V2 (\cref{ch:v2}) replaces the literature confinement laws by the first-principles values of this work. Its
pre-registered held-out test on the 6.3~nm film fails again: the threshold is too low by 0.23~V with the mobility
normalised to the on-current, and by 0.37~V with the pre-registered mobility rule.
\end{statusbox}

\section{Demonstrated conclusions}\label{sec:concl-demonstrated}

\begin{enumerate}
  \item The measured pattern is 2 $\approx$ 6.3~nm $\neq$ 13.2~nm in every metric (\evDATA). The values are $\Vthcc$ = 0.662/0.644/0.083~V, $\muFE$ = 12.15/10.95/50.17~\cmVs, $\SScc$ = 269.9/295.4/160.0~\mVdec, $\Ion$ = \SI{5.09e-7}{}/\SI{4.03e-7}{}/\SI{3.07e-6}{}~\Aum{} and off-floor \SI{4.6e-15}{}/\SI{1.5e-13}{}/\SI{6.6e-12}{}~\Aum{} (\cref{sec:inp-metrics-measured}). Since $\muFE(6.3) < \muFE(2)$, no monotonic law can pass through all three points.
  \item The work function, the electron affinity, the fixed charge and the pure confinement shift act as one number per curve (\evNUM). An increase of the gate WF by $+0.1$~eV gives a rigid shift of $+0.1000$~V with unchanged SS (\run{0048}, \run{0049}). $\Qf$ is rigid within 0.6~mV over $10^{-14}$..\SI{3e-7}{}~\Aum{} (\run{0001}$\rightarrow$\run{0003}, \run{0002}$\rightarrow$\run{0004}, \run{0005}$\rightarrow$\run{0007}), and $\dEc(2~\mathrm{nm}) = 0.353$~eV is equivalent to \SI{1.97e12}{cm^{-2}}. The complete confinement treatment ($\dEc$ together with $\mstar\rightarrow\Nc$) is only nearly rigid: the shift is 0.345~V at $10^{-14}$ and 0.302~V at $10^{-8}$~\Aum, and $\SScc$ is 289.1 against 303.7~\mVdec{} (\run{0012} against \run{0016}). Its non-rigid part is nearly reproduced by $\mstar\times0.8$ or by $\Nt\times1.09$ (log-linear scaling of \run{0050} and \run{0023}).
  \item The thresholds do not discriminate for or against confinement (\evNUM{} on \evDATA). A single shared offset fits with an rms of 0.1363~V with the confinement law and 0.1325~V without it; the leave-one-out values are 0.2045 and 0.1988~V, and the better variant changes with the choice of anchor film. Each reading requires one film-specific charge. With confinement, the 6.3~nm film needs \SI{1.64e12}{cm^{-2}} less positive charge (0.294~V); without confinement, the 13.2~nm film needs \SI{1.39e12}{cm^{-2}} more (0.248~V).
  \item The shared-parameter model misses the 6.3~nm threshold by $-0.294$~V (\run{0014}; \evOOS, failed). Because $\muband = 12.4$ is the V0 fit of the same curve, this is a test of the threshold only.
  \item No tested V1 variant reproduces the 6.3~nm on-state without degrading the subthreshold region (\evNUM). Eight variants of the charge and electrostatics (\run{0014}, \run{0015}, \run{0030}--\run{0035}) give a turn-on gap $\Vthlin-\Vthcc$ of 0.72--0.88~V, compared with 1.176~V measured. Variant C1 (\run{0052}) gives 0.939~V, and variant B0 (\run{0037}) gives 1.202~V but with $\SScc$ of 410.9 against 295.4~\mVdec. Smaller misses of the same sign exist at 2~nm (gap $-0.145$~V, $\gm$ $-8.6$\,\%) and at 13.2~nm ($-0.021$~V, $-2.6$\,\%).
  \item The constant-current threshold is confounded with the mobility (\evDATA/\evNUM). When the 13.2~nm curve is given the 2~nm mobility, $\Vthcc$ moves by $+0.108$ to $+0.115$~V, which is 19--20\,\% of the measured step of 0.579~V; the confinement term of the model realises 0.292~V (50.4\,\%).
  \item $\SSmin$ is not a valid metric (\evNUM). It moves by about 10~\mVdec{} between two curves with equal $\Ion$, so the claim ``removing confinement moves SS toward the data'' is an artefact (\cref{sec:der-metrics}).
  \item The calibrated state at DOS 384/192 is defined by \run{0038}, \run{0039} ($\times1.015074$) and \run{0040} ($\times1.020113$) (\evCAL). The measured and simulated values of $\Vthcc$ are 0.662/0.680, 0.644/0.653 and 0.083/0.054~V, and those of $\SSfix{10^{-11}}{10^{-10}}$ are 121.2/123.5, 124.5/124.5 and 91.8/90.9~\mVdec{} (13.2~nm floor-subtracted; raw value 136.1). The $10^{-10}$..$10^{-9}$ window is too soft by $+20$/$+6$/$+25$~\mVdec, and the fitted $\muband$ values are 17.69/11.58/61.60~\cmVs{} (\cref{sec:cal-final}).
  \item The 300~K transfer numerics converge where this was checked (\evNUM). At 2~nm, refinement of the DOS changes $\Ion$ by $+1.98$\,\% and then by $+0.48$\,\% (order 2.03, Richardson residual $+0.16$\,\%). A mesh refinement by $\times0.7$ passes at 6.3~nm (\run{0036}) and at 13.2~nm (\run{0028}), and $\Id$ is linear in $\mu$ to within $6\times10^{-6}$ (\cref{sec:num-convergence}).
  \item The elevated-temperature runs used the silicon default tmu = 1.5 (\evNUM). For this reason the predictions are registered in two exact mobility variants (\cref{sec:par-temperature}).
  \item \run{0027} is not an overlap test. Its malformed mesh collapsed the drain, and the run is retracted (\cref{sec:cal-param-control}).
  \item The mobilities published in the paper are obtained by applying the saturation formula to these curves (\evDATA, provenance). The formula gives 5.09 and 27.40~\cmVs, compared with the published 5.1 and 27.4, and it predicts 3.86~\cmVs{} at 6.3~nm (\cref{sec:der-mufe}).
  \item The computed gap opening of 1 and 2~nm \InO{} films lies 84--94\,\% in the conduction band (\evTHEORY, PBE; \cref{sec:dft-partition}); V1 had assumed 100\,\%. The band edges of the 2~nm film are bulk-like, not surface states. With these laws (model V2) the calibrated curves move rigidly by $-75$, $-24$ and $-10$~mV at 2, 6.3 and 13.2~nm. Re-calibrated with one shared $\Qf$ (\SI{1.383e12}{cm^{-2}}), V2 fits the anchors with an active RMSE of 0.065 and 0.097~dec (\run{0058}, \run{0059}; \evCAL).
  \item The pre-registered held-out test of V2 on the 6.3~nm film fails (\run{0060}; \evOOS, failed). $\Vthcc$ is $-0.23$~V off with the mobility normalised to $\Ion$ and $-0.37$~V off with the pre-registered mobility rule; $\Vthlin$ is $-0.58$~V off in both. The first-principles inputs shrink the threshold miss by about 0.04~V only. ATLAS agreed with the pre-registered (corrected) values within 0.3~mV in $\Vthcc$ (\cref{sec:v2-test}).
\end{enumerate}

\section{Strongly supported interpretations}\label{sec:concl-supported}

The first interpretation (S1) is that, within the calibrated model, the rise of the apparent mobility by a factor of $\approx$4.6 between 6.3 and 13.2~nm is a change in channel transport. With respect to the contacts, the net mobility-degradation coefficient is negative in all films, $\gm$ still rises at 3~V at 2~nm, and a thickness-independent $\Rsd$ would require $\rho_c\approx0.26~\Omega\,\mathrm{cm}^2$. With respect to electrostatics, campaign A changes $\gm$ by $\le2$\,\%. The quantum capacitance changes $C_{\mathrm{eff}}/\Cox$ by $\le1$\,\% at 2~nm and by 3--5\,\% at 6.3/13.2~nm (\evSURR). The roughness factor is $\times1.05$, and the tail partition changes from 0.791 to 0.826 (\evSURR). In total, 109\,\% of $\ln(\times4.58)$ is located in the fitted $\muband$. A limitation is that the bounds on $\Rsd$ belong to the calibration class ($\lesssim5\times10^4~\Omega\,\mu$m at 2~nm and $\lesssim1$--$2\times10^4~\Omega\,\mu$m at 13.2~nm). The details are given in \cref{sec:mech-transport}.

The second interpretation (S2) is that subthreshold conduction in all films is governed by the filling of an exponential acceptor tail near $\Ec$. One shared $\WTA$ (40~meV) reproduces the slope in the $10^{-11}$..$10^{-10}$~\Aum{} range within 2.3~\mVdec. ATLAS runs with ideal contacts reproduce the supralinear exponent $\alpha_H$ (1.59/2.14/1.29) and the negative $\theta$ through trap filling alone, and a DOS inversion validated on ATLAS curves gives measured/ATLAS ratios of 0.81--1.46 near $\Ec$ (conditional on the mobility). The specific DOS form of V1 is not supported, however: the deeper states are underestimated by $\times2.2$--2.8, and the $10^{-10}$..$10^{-9}$ window is too soft. The details are given in \cref{sec:mech-traps}.

The third interpretation (S3) is that the contacts do not create the thickness trend. The required $\Rsd$ would take up 83\,\% of $\Vd$. At fixed $\rho_c$ the contact fraction scales as $R_{\mathrm{sh}}^{-1/2}$, so a thickness-independent contact would penalise the thick film, and $\theta_{\mathrm{net}}$ has the wrong sign in all films (\evCORR{} + \evLIT). The details are given in \cref{sec:mech-contacts}.

The fourth interpretation (S4) is that the off-floors are caused neither by gate leakage (6.3 and 13.2~nm) nor by thermal generation. The 13.2~nm floor is flat within $\times1.06$ while $V_{gd}$ changes from 1.2 to 3.7~V, and the floor rises by $\times1432$ ($\propto t^{3.85}$) on an identical gate stack. Thermal generation falls short by 13--17 decades. The origin of the floors is not determined.

The fifth interpretation (S5) is that the published mobilities are values of the saturation formula taken outside the regime in which the formula applies. The exponent of $\Nt(t)$ in V1 descends from an analysis that was built on these values.

\section{Plausible hypotheses}\label{sec:concl-plausible}

The hypotheses that are plausible but not verified are collected in \cref{tab:concl-plausible}, together with the falsification criteria that were declared in advance.

\begin{table}[tbp]
  \centering
  \footnotesize
  \caption{Plausible but unverified hypotheses with their pre-declared falsifiers (\file{SYNTHESIS.md} \S C).}
  \label{tab:concl-plausible}
  \begin{tabularx}{\linewidth}{>{\raggedright\arraybackslash}p{0.04\linewidth}>{\raggedright\arraybackslash}p{0.28\linewidth}>{\raggedright\arraybackslash}X>{\raggedright\arraybackslash}X}
    \toprule
    \# & hypothesis & for / against & falsified if \\
    \midrule
    C1 & Material transition between 6.3 and 13.2~nm (crystallinity, grain coalescence, percolation) raising band mobility and effective positive charge together; \emph{working hypothesis} & co-located changes of $\Vth$, $\muFE$, $\SScc$ and floor (\evCORR); a 39~meV barrier difference suffices; ultrathin films lack long-range crystallinity \citep[p.~2]{Januar2026}. Against: no structural data at 6.3 or 13.2~nm; $N = 1$ & GIXRD/TEM show the same phase and grain size; or Hall mobility at 6.3~nm within 30\,\% of 13.2~nm; or replicate mean $\muFE(6.3)$ within 30\,\% of $\muFE(13.2)$ \\
    C2 & Donor or positive-charge step at 13.2~nm ($\approx$\SI{2.5e17}{}$\rightarrow$\SI{1e18}{cm^{-3}}) & needed by the no-confinement reading ($+$\SI{1.39e12}{cm^{-2}}); degenerate with confinement & C-V $V_{FB}(t)$ flat within 50~mV \\
    C3 & Confinement 0.26--0.30~V (0.14--0.38) at 2~nm & \evTHEORY{} (S3 Schrodinger-Poisson + PBE); this work's DFT gives $\dEc(1.98~\mathrm{nm}) = 0.281$~eV (PBE; $\times$1.06--1.11 with HSE06, \cref{sec:dft-partition}); not identifiable from the curves & optical-gap shift 2 vs 13.2~nm $<0.1$~eV \\
    C4 & 6.3~nm rigid offset is device-specific (spread, sweep history, stress) & PBS-extrapolated sweep shift 0.003--0.24~V & replicates give $\sigma(\Vth)<0.1$~V with mean $\ge0.2$~V off the shared law \\
    C5 & On-state shape needs a carrier-density-dependent (percolation-type) mobility & trade-off line of 3.2--8.5~mV gap per \mVdec{} of $\SScc$; untested & density-dependent mobility calibrated on 2 and 13.2~nm leaves the 6.3~nm gap miss $>0.1$~V at $\SScc\le330$~\mVdec \\
    C6 & Surface + bulk trap DOS ($D_s\approx$\SI{8.5e12}{cm^{-2}eV^{-1}} + $g_b\approx$\SI{1e19}{cm^{-3}eV^{-1}}) & ATLAS-validated inversion; energy axis shifts by $\kT\ln$(mobility ratio) & multi-frequency C-V per thickness \\
    C7 & Near-$\Ec$ states as partial contributor & \run{0052} recovers 33\,\% of the gap at $\SScc$ $+15.6$ & C1 at 768 levels loses the gain \\
    C8 & Floors are an ungated IWO path $\propto t^{3.85}$ & with the model's own donor and trap values, a strip carrying the V1 interface acceptors on one surface scales as $t^{4.24}$ and, with one factor fitted on 13.2~nm, gives $\times0.46$ and $\times0.64$ of the 2 and 6.3~nm floors (\cref{sec:v2-floor}); the 2~nm value is sensitive to two weakly known inputs; location not determined & $I_G\approx\Id$, or no $1/L$ scaling \\
    C9 & $\Qf$ is an interfacial ionized oxygen-vacancy layer & electrostatically possible; $\Qf$ absorbs $\pm0.2$~eV alignment & no C-V dispersion, no NBS response \\
    C10 & MTR temperature behaviour & registered P-MTR variant; untested & $\Delta\Ion$(2~nm, 358~K) $>+5$\,\% or $<-15$\,\% \\
    \bottomrule
  \end{tabularx}
\end{table}

\section{Rejected explanations}\label{sec:concl-rejected}

The rejected explanations and the quantitative reasons for their rejection are listed in \cref{tab:concl-rejected}. The term ``model-conditional'' means that the rejection holds within the DOS and electrostatics of V1.

\begin{table}[tbp]
  \centering
  \footnotesize
  \caption{Rejected explanations and the quantitative inconsistency (\file{SYNTHESIS.md} \S D).}
  \label{tab:concl-rejected}
  \begin{tabularx}{\linewidth}{>{\raggedright\arraybackslash}p{0.27\linewidth}>{\raggedright\arraybackslash}X>{\raggedright\arraybackslash}p{0.16\linewidth}}
    \toprule
    explanation & quantitative reason & class \\
    \midrule
    6.3~nm thickness error (A1, \run{0031}) & $-1$~nm gives $+0.036$ of the required $+0.294$~V (12\,\%); closing via $\dEc$ needs $t\approx2.1$~nm & \evNUM, model-cond. \\
    Fewer donors (A2, \run{0032}) & $+0.030$~V, $\times10$ short & \evNUM \\
    Front $\Dit$ (A4, \run{0033}) & $\times5$ gives $+0.051$~V, gap unchanged; needs $\times23$--24 at 6.3~nm only & \evNUM \\
    Back-surface charge alone (A3, \run{0030}) & matches $\Vthcc$ but $\SScc$ off by $-64.1$~\mVdec; offset at $10^{-7}$ worsens to $+0.285$~V & \evNUM, model-cond. \\
    Combination A5 (\run{0035}) & 0.215 of 0.294~V with $\SScc$ $-45.9$~\mVdec & \evNUM \\
    S2 re-partition (B0, \run{0037}) & $\SScc$ 410.9 vs 295.4; $\SSfix{10^{-11}}{10^{-10}}$ $+48.8$~\mVdec & \evOOS, failed \\
    Near-$\Ec$ band as full explanation (C1, \run{0052}) & gap $+0.119$ vs $+0.28$~V predicted; $\Ion$ $+16.0$\,\% & \evOOS, partly failed \\
    Any rigid charge as whole 6.3~nm miss & offset varies $-0.007$ to $+0.141$~V between $10^{-9}$ and \SI{2e-7}{}~\Aum{} (\run{0015}) & \evNUM{} on \evDATA \\
    Thickness-independent $\Rsd$ as the trend & needs $\Rsd W\approx1.15\times10^6~\Omega\,\mu$m (83\,\% of $\Vd$); $\theta_{\mathrm{net}}$ wrong sign & \evCORR{} + \evLIT \\
    Confinement-raised Pd/IWO barrier as the step & needs $\Delta\phi\ge0.20$~eV; $\dEc(6.3) = 0.07$~eV & \evCORR \\
    Roughness factor as mobility step & $\times1.051$ vs $\times4.58$ & \evCAL \\
    Thickness-fluctuation scattering ($\propto t^6$) & implies $\mu(2)$ = 0.017~\cmVs{} ($\times700$ off) & \evSURR \\
    Back-surface Coulomb scattering & $\times2700$ short at 6.3~nm & \evSURR \\
    Tail partition as mobility step & P changes $+4$\,\%; needs $\Nt\times6.7\times10^3$ & \evSURR \\
    Uniform donors as 6.3$\rightarrow$13.2 $\Vth$ step & steepest slope ratio 2.35 vs 9--20 measured & \evSURR{} + \evDATA \\
    Gate leakage as floor (6.3, 13.2~nm) & $\le\times1.06$ change while $V_{gd}$ 1.2$\rightarrow$3.7~V & \evDATA \\
    Thermal generation as floor & 13--17 decades short & calculation \\
    Power laws in $t$ for $\muFE$, $\Ion$, $\Vth$ & non-monotonic data; LOO miss $\times2.63$ ($\muFE$), $\times3.76$ ($\Ion$) & \evDATA \\
    $\dEc\propto t^{-1.38}$ beyond $\approx3$~nm as physics & exceeds the infinite-well bound above 2.98~nm ($\times1.59$ / $\times2.51$ at 6.3 / 13.2~nm) & \evTHEORY \\
    ``Removing confinement improves SS'' & window artefact (\cref{sec:der-metrics}); fixed-current SS within 1--2~\mVdec & \evNUM \\
    ``Thresholds require/refute confinement'' & rms 0.136 vs 0.133~V & \evNUM{} on \evDATA \\
    ``Contact geometry irrelevant'' & never tested (\run{0027} malformed) & audit \\
    \bottomrule
  \end{tabularx}
\end{table}

Several explanations are neither supported nor rejected and remain not determined. These are hysteresis, bias stress during the sweep, device-to-device spread, the measurement ambient, and a thickness error beyond the model-conditional bound.

\section{Publication-ready claims}\label{sec:concl-claims}

The following four statements are the strongest claims that the present study supports (\file{SYNTHESIS.md} \S J). They are formulated so that they can be quoted without change.

\begin{enumerate}
  \item \emph{(\evNUM; demonstrated.)} On a single transfer curve per thickness, the confinement-induced conduction-band shift, the gate work function, the channel electron affinity and a fixed interface charge enter as one rigid threshold offset (reproduced to within 1~mV in drift-diffusion simulations). A confinement shift of about 0.3~eV at 2~nm, expected from effective-mass estimates anchored on first-principles slab calculations, is therefore an assumed input, not a result, and the measured thresholds are described equally well with or without it (single-offset rms 0.136 vs 0.133~V).
  \item \emph{(\evDATA{} + \evNUM; demonstrated for these devices.)} The 2 and 6.3~nm devices have similar thresholds and field-effect mobilities, whereas the 13.2~nm device shows a 0.56--0.58~V lower threshold and a 4.1--4.6-fold higher apparent mobility. With the calibrated tail-state and donor laws and a thickness-independent fixed charge, a single offset does not reproduce all three thresholds within 0.1~V, with or without the confinement term (largest residuals 0.19 and 0.18~V); at least one further thickness-dependent quantity is required, which the present data do not identify. About 0.11~V of the threshold difference follows from the constant-current threshold definition combined with the mobility difference.
  \item \emph{(Strongly supported; model-conditional.)} Within the calibrated trap-limited model, the 4.6-fold increase in apparent mobility between 6.3 and 13.2~nm cannot be attributed to source/drain resistance, interface roughness, electrostatics or tail-state filling; its physical origin, for example a thickness-dependent microstructure, is not determined by the present data.
  \item \emph{(\evNUM{} on \evDATA; demonstrated negative result.)} The 6.3~nm device departs from the shared calibration in two separable ways: a rigid threshold offset (0.29~V, equivalent to $1.6\times10^{12}$~\cmsq{} of fixed charge) and a gate-voltage-dependent on-state shape (turn-on delay 0.31--0.36~V too small, transconductance 17--23\,\% too low). No tested change of fixed charge, donor density, interface, tail-state or near-band-edge acceptor density, film thickness or back-surface charge reproduces the shape without degrading the subthreshold slope. With one device per thickness, the offset cannot be distinguished from device-to-device variation, and the shape remains unexplained; a limitation of the constant-mobility model is the leading untested candidate.
\end{enumerate}

\section{Corrections to earlier package documents}\label{sec:concl-corrections}

None of the earlier documents has been edited. The corrections in \cref{tab:concl-corrections} must be recorded separately (\file{SYNTHESIS.md} \S K), and they are applied throughout this report.

\begin{table}[tbp]
  \centering
  \footnotesize
  \caption{Principal corrections to earlier package documents.}
  \label{tab:concl-corrections}
  \begin{tabularx}{\linewidth}{>{\raggedright\arraybackslash}p{0.17\linewidth}>{\raggedright\arraybackslash}p{0.3\linewidth}>{\raggedright\arraybackslash}X}
    \toprule
    document & superseded statement & corrected wording \\
    \midrule
    FINAL\_STATUS & ``15 real ATLAS launches'' & 52 recorded launches (\run{0001}--\run{0052}), of which 49 produced scored currents; \run{0010} aborted, \run{0011} stopped, \run{0018} timed out; \run{0027} malformed and retracted \\
    FINAL\_STATUS & ``The final set is four runs'' & final calibrated set \run{0038}--\run{0040} (DOS 384/192; $\muband$ 17.69/11.58/61.60); \run{0014} is the shared-parameter threshold test; \run{0012}--\run{0015} superseded \\
    FINAL\_STATUS & $\SSmin$ column; row 0014 ``VALIDATION'' & fixed-current SS, $\SScc$, gap, $\gm$; ``shared-parameter out-of-sample threshold test: failed, $-0.294$~V'' \\
    FINAL\_STATUS & ``ONE fitted electrostatic value'' & fitted: shared $\Qf$, 6.3~nm $\Qf$, $\muband\times3$, $\Nt$ anchor and $\WTA$, $\Nd$-law constants; assumed: WF, $\chi$, $\Dit$, deep Gaussian, $\varepsilon$(\AlO) \\
    FINAL\_STATUS & 6.3~nm ``reproduced in SHAPE'' & fixed-current SS reproduced ($-2.2$~\mVdec); on-state not (gap $-0.314$~V, $\gm$ $-16.6$\,\%) \\
    FINAL\_STATUS & Conclusion 1: confinement ``accounts for'' the 0.58~V & thresholds do not discriminate (rms 0.136 vs 0.133~V); realised term 0.29~V (50\,\%); confinement is an assumed input \\
    FINAL\_STATUS & ``publication-ready as a calibrated parameter extraction'' & a calibrated, explicitly non-unique description of three single devices; not validated; submission requires Tiers 0--3 (\cref{sec:val-checklist}) \\
    SENSITIVITY\_RESULTS & overlap\_4um (\run{0027}) ``MEASURED'' & RETRACTED: malformed deck \\
    SENSITIVITY\_RESULTS & WF, $\mstar$, 6.3~nm ``NOT RUN'' & superseded by \run{0048}--\run{0051} and \run{0030}--\run{0036} \\
    NUMERICAL\_CONV. & V0 checks apply to V1 decks & carried over by lineage; not re-demonstrated on the V1 2~nm deck \\
    NUMERICAL\_CONV. & DOS C' ``FAIL'' & resolved by recalibration at 384/192; order established only at 2~nm \\
    LIMITATIONS & \#5 ``$\approx$12 $\rightarrow$ $\approx$57'' & 11.58 $\rightarrow$ 61.60~\cmVs{} (DOS 384/192) \\
    LIMITATIONS & \#11 overlap insensitive & contact geometry and resistance untested \\
    SUPERVISOR\_SUMMARY & ``all 14 launches''; ``confinement story holds ... it was tested'' & 52 recorded launches, of which 49 produced scored currents; FINAL\_STATUS Conclusion 1 wording \\
    EVIDENCE\_BRIEF & 13.2~nm no-confinement $\Qf\approx$\SI{1.86e12}{} & \SI{1.434e12}{cm^{-2}} \\
    S2 & re-partition ``strongly supported''; ``linear regime'' & rejected (B0); regime claim withdrawn; 31.8~nm ``$\muFE$'' is not a mobility \\
    S3 / S1 / S6 & ``mildly disfavour''; A6 ``favours no-QC''; P2, A6 pre-registered & ``do not discriminate''; ``neutral''; POST \\
    S4 & 6.3~nm floor $\pm3$\,\% & $\pm45$\,\% \\
    PREDICTIONS\_REGISTER & prose multipliers 1.19665/1.30426 & exact 1.19669/1.30441 (tables already exact); addendum only \\
    \bottomrule
  \end{tabularx}
\end{table}

\section{Outlook}\label{sec:concl-outlook}

The study has turned an apparently complete thickness-resolved TCAD fit into a map of what three single transfer curves can and cannot identify. The next scientific steps follow the order of dependencies given in \cref{sec:val-checklist}. First, the registered predictions must receive an external timestamp and must then be compared, with the unchanged extractor and without refitting, against the existing 2~nm temperature data in the supplement of the target paper. This comparison is the first genuine validation test, and it decides between the P-MTR and P-phonon transport variants. Second, replicate devices (3--5 per thickness, measured with dual sweeps) are needed before any film-specific statement about the 6.3~nm device can be made. Third, thickness metrology, the C-V flat-band voltage as a function of thickness, TLM or L-series contact measurements and a temperature series on all three films would remove the principal degeneracies, namely those between work function, affinity, fixed charge and confinement, between band mobility and tail density, and between $\Rsd$ and the tail exponent. Structural characterisation (GIXRD, TEM, grain size) at 6.3 and 13.2~nm is the direct test of the working hypothesis C1. A carrier-density-dependent mobility model, first-principles calculations for W-doped \InO{} and further thicknesses belong to a second study (\cref{sec:val-second-paper}).

\begin{keyresult}[title={Key result: conclusions}]
The study demonstrates that the measured devices follow the pattern 2 $\approx$ 6.3 $\neq$ 13.2~nm, that confinement, WF, $\chi$ and $\Qf$ act as one rigid number per curve, and that the thresholds are neutral with respect to confinement (rms 0.136 vs 0.133~V). The 6.3~nm miss decomposes into a rigid part of 0.294~V and an on-state shape part that no tested V1 change reproduces and that remains unexplained. It is strongly supported, although conditional on the model, that the rise of the mobility by a factor of $\approx$4.6 is a change in transport within the calibrated model that is not caused by contacts, roughness, electrostatics or tail filling, and that subthreshold conduction is governed by tail filling. A material transition between 6.3 and 13.2~nm remains the working hypothesis; it is plausible but untested. The model is calibrated and non-unique but not validated, and the registered predictions (\run{0041}--\run{0047}) await independent data. Replacing the literature confinement laws by first-principles values of this work (model V2) changes the calibrated curves by at most 75~mV and does not repair the held-out 6.3~nm prediction, which fails its pre-registered test again ($-0.23$~V).
\end{keyresult}
